% =============================================================
% LSVJ-S: Synthesize, Gate, Verify
% Runtime Proof-Obligation Synthesis for Trustworthy Agent Decisions
% NeurIPS 2024 style, [preprint] option for arXiv
%
% DRAFT 2026-04-21. Numerical results in Section 4 are marked
% [TBD] until R002 pilot and W6-W7 runs complete. Do NOT submit
% without replacing placeholders and running compile + auto-review.
% =============================================================
\documentclass{article}
\usepackage[preprint]{neurips_2024}

\usepackage{amsmath,amssymb,amsfonts}
\usepackage{amsthm}
\usepackage{algorithmic}
\usepackage{graphicx}
\usepackage{textcomp}
\usepackage{xcolor}
\usepackage{booktabs}
\usepackage{multirow}
\usepackage{hyperref}
\usepackage{listings}
\usepackage{mathtools}
\usepackage{natbib}
\usepackage{enumitem}

\theoremstyle{definition}
\newtheorem{definition}{Definition}
\newtheorem{proposition}{Proposition}

\hypersetup{colorlinks=true,linkcolor=blue,citecolor=blue,urlcolor=blue}

\lstset{basicstyle=\ttfamily\small,breaklines=true,frame=single,captionpos=b,
        keywordstyle=\bfseries,commentstyle=\itshape\color{gray!70}}
\lstdefinelanguage{cozo}{
  keywords={not,and,or,true,false,discharged},
  sensitive=true,
  comment=[l]{\%},
  morestring=[b]"
}

\newcommand{\todo}[1]{\textcolor{red}{\textbf{[TBD: #1]}}}
\newcommand{\placeholder}[1]{\textcolor{red}{\textbf{[#1]}}}

\title{LSVJ-S: Synthesize, Gate, Verify --- Runtime Proof-Obligation
Synthesis for Trustworthy Agent Decisions}

\author{%
  Dongcheng Zhang\thanks{Correspondence: \texttt{zdclink@gmail.com}.
  Work performed while at BlueFocus Communication Group.} \\
  Independent \\
  Beijing, China \\
  \texttt{zdclink@gmail.com}
  \And
  Yiqing Jiang \\
  Tongji University \\
  Shanghai, China \\
  \texttt{jyq.russel@gmail.com}
}

\begin{document}
\maketitle

% ================================================================
\begin{abstract}
LLM-based agent runtime authorization suffers a fundamental epistemic
gap: the LLM's ``allow'' verdict is an opaque label with no mechanical
connection to the facts it claims to reason over. Existing hybrid
architectures stack symbolic layers alongside the LLM, but the LLM's
decision is never verified \emph{by} the symbolic substrate. We present
LSVJ-S, a Synthesis--Proof-Obligation (SPO) protocol in which the LLM
synthesizes a per-decision Datalog rule as its committed proof of trust,
and a 4-stage compile-time gate---(b.1) parse, (b.2) type-check against
a typed owner-harm primitive schema (Class A/B/C), (b.3) syntactic
non-triviality, and (b.4) the compound \emph{perturbation-sensitive AND
has-decisive-primitive} check---admits the rule only if all stages pass.
Gate b.4 is grounded in Self-Consistency~\citep{wang2022selfconsistency}
applied to LLM-synthesized Datalog; to our knowledge this is the first
structured Self-Consistency application to runtime safety-rule
synthesis. On a hard-negative held-out slice (v3-HN), LSVJ-S is compared
against a 12-rule Static Expert-Authored Policy (SEAP) baseline that is
structurally isomorphic to the PCAS offline-compilation
model~\citep{palumbo2026pcas}. \placeholder{Headline result TBD after
W6 run: report LSVJ-S TPR/FPR with Wilson 95\% CI and the
paired-McNemar per-primitive causal attribution from B4.} We explicitly
acknowledge the Self-Justification Trap---that a misaligned proposer
can synthesize a collusion rule---and provide empirical defense via
adversarial synthesis stress testing (B6). Code, ontology schema, and
v3-HN protocol are released at \url{https://github.com/dario-github/nous}.
\end{abstract}

% ================================================================
\section{Introduction}
\label{sec:intro}

Autonomous LLM agents are increasingly authorized to invoke
consequential tools---sending emails, executing shell commands,
modifying production data---based on a single verdict label emitted by
an LLM-based safety gate. The verdict is opaque: an ``allow'' stands
on the LLM's internal reasoning but is not mechanically connected to
the system state it claims to reason over. Our companion work on
Owner-Harm~\citep{zhang2026ownerharm} documents the operational
consequence: a compositional gate that reaches 100\% TPR on generic
criminal-harm benchmarks collapses to 14.8\% TPR (95\% CI: 5.9--32.5\%)
on the AgentDojo injection slice, where the harm is mediated through
prompt injection into legitimate tool vocabulary.

Internal results on our owner-harm benchmark show a sharper pathology.
A compositional gate with a deterministic Datalog layer (L1) and a
triviality filter (L2) reaches zero training loss; yet on the held-out
slice the same L1+L2 subsystem achieves only \textbf{12.7\%} TPR. Most
of the held-out detection (+62.6\,pp) is carried by an LLM semantic
layer (L3) whose verdicts are, by construction, unverified against the
symbolic substrate. A \texttt{no\_kg} ablation on the training set
further shows that the knowledge graph's contribution is statistically
indistinguishable from zero ($\Delta L = 0$), suggesting the L3
layer's held-out gains are not symbolically grounded but rather a
function of the LLM's opaque pattern matching.

This paper asks: \emph{can the symbolic substrate become the verifier,
rather than a stage adjacent to the LLM decision?}

\paragraph{Why existing approaches fall short.}
Three families of runtime agent safety enforcement dominate the
literature (cf.\ Section~\ref{sec:related}):
\begin{itemize}[leftmargin=1.5em]
\item \textbf{LLM-only classifiers}
(Llama Guard~\citep{inan2023llamaguard};
ShieldGemma~\citep{zeng2024shieldgemma}): opaque labels without an audit
artifact at the symbolic layer.
\item \textbf{Hybrid stacking} (Nous~\citep{zhang2026ownerharm};
GuardAgent~\citep{xiang2024guardagent};
TrustAgent~\citep{hua2024trustagent};
ShieldAgent~\citep{chen2025shieldagent}): the LLM and the symbolic
engine run in parallel, not in a verifier--proposer relationship.
\item \textbf{Offline policy compilers}
(PCAS~\citep{palumbo2026pcas}; Solver-Aided~\citep{roy2026solveraided};
AgentSpec~\citep{wang2026agentspec}): strong deterministic enforcement
inside closed, pre-audited vocabularies; no mechanism for
open-vocabulary contextual cases the policy author did not anticipate.
\end{itemize}
Each family forfeits either symbolic grounding (families 1--2) or
open-vocabulary reach (family 3). A verifier--proposer architecture
where the LLM \emph{commits} to a per-decision symbolic proof
obligation, checked at run time, remains an unfilled gap.

\paragraph{Insight.}
Instead of producing an opaque ``allow'' label, the LLM should produce
a \emph{committed proof obligation}---a typed Datalog rule over a
bounded primitive vocabulary---and ``allow'' stands only when (a) the
rule passes a 4-stage compile-time soundness gate and (b) the rule
discharges against live facts from the agent's knowledge graph and
host-function observables. The gate's authority does not depend on
LLM capability: it rejects unconditionally on any stage failure,
fail-closed to \texttt{confirm}. This turns the symbolic substrate
into a verifier, not a witness.

Explicit scope: LSVJ-S is not a formal proof system. The gate is a
\emph{soundness check} that rejects invariant and trivial rules, not a
completeness guarantee that admits all and only correct rules. The
contribution is protocol-level plus empirical.

\paragraph{Contributions.}
\begin{enumerate}[leftmargin=1.5em]
\item \textbf{SPO protocol.} Per-decision Datalog rule synthesis over
a typed owner-harm primitive schema partitioned into Class A
(KG-grounded), B (deterministic host function), and C (sealed
sub-oracle); ``allow'' requires both gate admission and live-fact
discharge.
\item \textbf{4-stage compile-time gate} including the compound b.4
check (perturbation-sensitive AND has-decisive-primitive). Gate b.4 is
structured Self-Consistency~\citep{wang2022selfconsistency} applied to
LLM-synthesized Datalog, addressing the Self-Justification Trap.
\item \textbf{Empirical evaluation on v3-HN:} (a) a 5-configuration
main comparison with Wilson 95\% CI, (b) 9 Tier-1 ablation conditions,
(c) per-primitive paired-McNemar causal attribution, and (d) an
adversarial-synthesis stress test (B6) that empirically bounds the
Self-Justification Trap defense.
\end{enumerate}

% ================================================================
\section{Related Work}
\label{sec:related}

\subsection{Three Families of Runtime Agent Safety Enforcement}

\paragraph{Family 1 --- Static DSL + deterministic reference monitor.}
PCAS~\citep{palumbo2026pcas}, AgentSpec~\citep{wang2026agentspec}, and
ClawGuard~\citep{clawguard2026} author rules offline, compile them,
and enforce deterministically. PCAS is our closest prior art and the
structural model for our SEAP baseline. Following a clean metaphor:
PCAS is a \emph{statute} model---laws written before the crime,
deterministically enforced---while LSVJ-S is a \emph{case-law} model
in which the LLM argues the legality of its specific action at the
moment of execution, and a compile-time soundness gate checks the
non-triviality of the argument before admitting. The architectural
delta is the axis of rule authorship (offline human-audited vs.\
runtime LLM-synthesized) and the resulting need for a compile-time
gate, which PCAS does not require because its rules are pre-audited.
Scope qualification: PCAS-style offline policies dominate LSVJ-S when
the domain is closed and pre-enumerable; LSVJ-S targets the
complement---open-vocabulary contextual edge cases.

\paragraph{Family 2 --- Formal-logic constraint verification.}
Solver-Aided~\citep{roy2026solveraided} translates natural-language
policies to SMT-LIB offline; the LLM does not synthesize the
constraint. Agent-C~\citep{agentc2026} interleaves SMT with
constrained decoding at token-generation time; that is a
model-internal layer whereas LSVJ-S is a model-agnostic wrapper. The
constraint-source axis is the delta: static pre-translated vs.\
per-decision LLM-synthesized.

\paragraph{Family 3 --- Probabilistic and learned guardrails.}
Pro2Guard~\citep{he2025pro2guard} learns a DTMC from traces;
ShieldAgent~\citep{chen2025shieldagent} uses probabilistic rule
circuits; GuardAgent~\citep{xiang2024guardagent} and
TrustAgent~\citep{hua2024trustagent} rely on LLM-based agent-of-agent
structures. LSVJ-S differs: classical non-probabilistic Datalog
synthesis plus a compile-time soundness gate, no probabilistic scoring
inside the verifier.

\subsection{Foundational Methods LSVJ-S Builds On}

\paragraph{Self-Consistency.}
Gate b.4's perturbation-sensitivity check is structurally
Self-Consistency~\citep{wang2022selfconsistency} applied to
LLM-synthesized Datalog rules: the rule is evaluated on live bindings
and on $N{=}5$ perturbed bindings, and is rejected if the outcome is
invariant across all trials. We credit this lineage explicitly: b.4
is \textbf{not} an independent invention relative to Self-Consistency.
LSVJ-S's novelty is the \emph{combination} of structured-output
Datalog synthesis, perturbation-sensitivity as gate, and
decisive-primitive causal attribution, applied to the runtime agent
safety setting---none of the three sub-steps is independently novel.

\paragraph{Program synthesis with symbolic verification.}
Scallop~\citep{li2023scallop} introduced differentiable Datalog for
neurosymbolic learning; Lobster~\citep{lobster2026} extends the
neurosymbolic stack. Both are offline-training systems. LSVJ-S is
classical runtime Datalog with zero proposer training---this is a
deliberate design choice defended in Section~\ref{sec:method:tradeoffs}
on epistemic grounds, not only budgetary ones.

\paragraph{Formal safety for tool use.}
\citet{doshi2026verifiable} provide the formal-safety-for-tool-use
motivation via STPA applied to MCP; we cite as motivation only.

\subsection{SEAP: The PCAS Analog Baseline}
Our primary baseline, \emph{Static Expert-Authored Policy (SEAP)}, is
a pre-registered 12-rule Datalog menu over the same A/B/C primitive
vocabulary used by the LSVJ-S proposer (Appendix~A). SEAP is
structurally isomorphic to PCAS offline compilation: offline-authored,
no per-decision synthesis, same predicate universe. The main-table
comparison LSVJ-S vs.\ SEAP is therefore an apples-to-apples test of
the architectural axis this paper introduces: does per-decision
runtime synthesis add TPR beyond an expert-authored policy of the
same vocabulary?

% ================================================================
\section{Method}
\label{sec:method}

\subsection{Owner-Harm Primitive Schema (Class A/B/C)}
\label{sec:method:schema}
A typed primitive vocabulary is the contract between proposer and
verifier: it enables stage b.2 (type-check) without per-rule human
review. We partition six primitives across three classes (full schema
in Appendix~A):
\begin{itemize}[leftmargin=1.5em]
\item \textbf{Class A --- KG-grounded}: a Cozo Datalog query over the
live NousDB; returns the empty relation if the current KG state does
not support it. Examples:
\texttt{is\_inner\_circle(recipient\_id)},
\texttt{owner\_has\_directed(action\_id, is\_directed)}.
\item \textbf{Class B --- deterministic host function}: pure Python,
no LLM, returns \texttt{bool}. Examples:
\texttt{body\_length\_within\_budget(body, budget)},
\texttt{action\_matches\_allowlist(action\_id, list)}.
\item \textbf{Class C --- sealed sub-oracle}: a fresh session of the
same base model with schema-bound input and divergence-calibrated
output. Honest scope label: \emph{procedural} independence, not
representation-level (the proposer and the Class C sub-oracle share
weights and plausibly share failure modes). Different-model Class C
is listed as future work. Examples:
\texttt{body\_reveals\_inner\_relation(body, recipient\_id, p\_rev)},
\texttt{action\_is\_reversible(action\_id, reversible)}.
\end{itemize}
Each class admits a distinct b.4 perturbation operator
(A: reshuffle matching tuples within-type; B: negate bool return;
C: flip a pre-recorded sub-oracle output with no additional LLM call).
The 6-seed library and the 12-rule SEAP menu are prior development
and baselines, not contributions of this paper.

\subsection{Proposer: Grammar-Constrained Synthesis}
\label{sec:method:proposer}
The proposer is a frozen production LLM
(\placeholder{GLM-4.7-flashx in the R002 pilot; final paper reports
the model selected after W2 per Section~\ref{sec:method:tradeoffs}})
invoked with XGrammar~\citep{xgrammar2024}-constrained decoding over
a Lark EBNF subset of the Cozo rule grammar. The proposer emits a
single JSON object
\texttt{\{decision, synthesized\_obligation: \{rule\_body, decision\}\}}.
Grammar constraint enforces well-formedness before stage b.1,
reducing gate rejection cost and simplifying parse-failure analysis.

Each synthesized rule has the form
$\mathtt{?[discharged] := \text{body},\ discharged = \phi(body)}$
where $\text{body}$ is a conjunction of A/B/C primitive calls with
fixed arity, and $\phi$ is a Boolean expression over bindings
introduced by the body. Figure~\ref{fig:example-rule} shows a
representative rule for the \texttt{send\_email} tool invocation.

A six-rule seed library is provided as few-shot context to bootstrap
the grammar. To diagnose over-anchoring, we report the
\emph{verbatim-match rate}---the fraction of synthesized rules that
exactly match a seed rule or SEAP rule---and treat rates above 50\%
as a signal that synthesis collapsed to template completion.

\begin{figure}[t]
\centering
\begin{lstlisting}[language=cozo,label=lst:example]
?[discharged] :=
  is_inner_circle(recipient_id),                    % Class A
  body_reveals_inner_relation(body,                 % Class C (sealed)
                              recipient_id, p_rev),
  owner_has_directed(action_id, is_directed),       % Class A
  discharged = (not p_rev) or is_directed
\end{lstlisting}
\caption{Example LSVJ-S rule for a \texttt{send\_email} tool call.
Allow iff the recipient--body relation is unlikely to be revealing,
\emph{or} the action was owner-directed. Class labels annotate the
primitive's verification path.}
\label{fig:example-rule}
\end{figure}

\subsection{The Four-Stage Compile-Time Gate}
\label{sec:method:gate}
The gate fails closed to \texttt{confirm} on any stage failure; its
authority is independent of proposer capability.

\paragraph{b.1 --- Parse.}
The rule is fed to the Cozo parser. Parse failure yields
\texttt{confirm}. In the R002 pilot (100 cases sampled from v3
held-out, \texttt{glm-4.7-flashx} proposer, no grammar-constrained
decoding), we observed 97/100 responses captured (3 API failures)
and, on captured responses, \textbf{97.9\%} parse rate ($95/97$,
after carrying the outer \texttt{decision} field into the inner
\texttt{obligation.decision} when the LLM format-confused) and
\textbf{96.9\%} end-to-end gate-admission rate ($94/97$).
Per-category parse rate is $100\%$ on 7/8 categories; Unauthorized
Autonomy is the sole lagger at $87.5\%$ ($14/16$) due to two
residual LLM-emission bugs (one truncated JSON mid-string, one
unescaped backslash in a regex literal). These are proposer-emission
limits, not gate-logic failures; the gate correctly rejects them via
b.1. The original $\le 1\%$ failure target assumed grammar-constrained
decoding (e.g.\ XGrammar~\citep{xgrammar2024}), which is not available
via the \texttt{openai-completions} API on our deployment; reaching
the target will require a grammar-decoder integration (future work).

\paragraph{b.2 --- Type-check.}
Every primitive call must match a declared primitive with correct
arity and class. Unknown predicates, arity mismatches, or class
violations yield \texttt{confirm}. This is the structural payoff of
the A/B/C partition.

\paragraph{b.3 --- Syntactic non-triviality.}
The body must reference at least one Class A or Class B primitive,
and the head expression must not be a constant or literal Boolean.
This stage catches tautologies (e.g.\ \texttt{discharged = true})
that pass b.1 and b.2 but are semantically vacuous.

\paragraph{b.4 --- Perturbation-sensitive AND has-decisive-primitive.}
The compound stage comprises two logically independent checks sharing
evaluation machinery:
\begin{itemize}[leftmargin=1.5em]
\item \emph{Perturbation-sensitivity.} Run the rule on live bindings
to obtain $\mathtt{discharged\_real}$. Then run $N{=}5$ trials on
perturbed bindings (per-class operator from
\S\ref{sec:method:schema}). Reject the rule if the outcome is
identical across all $N$ trials. Calibration (W3 pilot,
\placeholder{40-rule hand-labeled reference set TBD}): false-rejection
$\le 10\%$ and false-admission on content-invariant rules $\le 2\%$.
\item \emph{Has-decisive-primitive.} For each primitive, flip its
binding and check whether the outcome changes. Reject the rule if no
primitive is decisive. This ensures every admitted rule has an
attributable cause and provides an audit artifact as a byproduct.
\end{itemize}
We disclose the honest scope: b.4 rejects invariant rules, not
weakly-dependent ones. A rule whose primitive flip rate lies in
$[0.10, 0.25]$ passes $N{=}5$ with probability $\approx 0.67$. This is
a feature of the zero-training protocol; a stronger gate would
require training the proposer, which would collapse verifier
independence.

\subsection{Decisive-Primitive Tracing (Diagnostic)}
\label{sec:method:tracing}
We define
$\mathtt{decisive\_primitives}(\mathit{rule}, \mathit{bindings})$ as
the set of primitives whose binding flip changes the discharge
outcome. This is a \emph{diagnostic} utility, not a gate component.
The set is recorded in
\texttt{proof\_trace.lsvj\_obligations[i].binding\_decisive\_primitive}
and used for (i) the paired-McNemar per-primitive causal attribution
in B4, (ii) case studies in B5, and (iii) a neg-control anti-gaming
check (fraction of LSVJ-S discharges on neg-control that show
\emph{any} Class~A primitive decisive must be $\le 20\%$;
pre-registered failure path: recalibrate seed library).

\subsection{Full Inference Path}
\label{sec:method:path}
Outside L3, the Nous pipeline is unchanged:
\begin{enumerate}[leftmargin=1.5em]
\item \textbf{L1} Datalog exact-match $\to$ exit on match
\item \textbf{L2} triviality filter $\to$ exit on trivial
\item \textbf{L3-LSVJ-S}: (a) proposer $\to$
\texttt{\{decision, synthesized\_obligation\}};
(b) 4-stage gate (\S\ref{sec:method:gate});
(c) execute rule over live bindings
(A: Cozo; B: host function; C: sealed sub-oracle);
(d) decisive-primitive tracing;
(e) verifier policy below.
\item \textbf{L4} hijacking filter, unchanged from
\citet{zhang2026ownerharm}.
\item Final verdict with
\texttt{proof\_trace.\{rule\_text, class\_usage, discharge\_result,
decisive\_primitives\}}.
\end{enumerate}

\paragraph{Verifier policy.}
If the LLM says \texttt{block} or \texttt{confirm}, accept as-is. If
the LLM says \texttt{allow} \emph{and} the rule is admitted \emph{and}
the rule discharges against live bindings, output \texttt{allow}.
Otherwise upgrade to \texttt{confirm}.

\subsection{Honest Design Tradeoffs}
\label{sec:method:tradeoffs}

\paragraph{Zero-training rationale.}
(i) Grammar-constrained decoding delivers schema adherence without
fine-tuning. (ii) Supervised fine-tuning on obligations would
co-adapt the proposer to the verifier's acceptance criterion,
collapsing verifier--proposer independence, which is the very
property on which the compile-time gate's authority rests.
(iii) The API-only budget (\$120) makes training infeasible. All
three constraints are honest; (ii) is the dominating epistemic
argument.

\paragraph{Self-Justification Trap.}
The proposer writes both the action (implicitly, via the LLM's
earlier verdict) and the admitting rule. A misaligned or
prompt-injected proposer can synthesize (i) a tautology, (ii) a rule
referencing only rarely-enforcing primitives, (iii) a collusion rule
that passes b.1--b.3 while remaining semantically vacuous, or
(iv) an OR-disjunction where only one term drives discharge. Stages
b.3 and b.4 are the primary defenses; B6
(\S\ref{sec:exp:b6}) empirically tests them against an attacker-mode
proposer.

\paragraph{Class C same-model limitation.}
The Class C sub-oracle is a fresh session of the same base model with
schema-bound input. This is procedural, not representational,
independence: shared weights plausibly share failure modes. Named
limitation; different-model Class C is future work.

\paragraph{Decorative-KG contingency.}
If the B4 per-primitive causal test shows null effect on all Tier-1
primitives, we pre-commit to a framing pivot from \emph{symbolic
substrate verifier} to \emph{JIT-policy gate}
(\S\ref{sec:disc:pivot}). This is pre-registered, not a post-hoc
save.

% ================================================================
\section{Experiments}
\label{sec:exp}

\subsection{Datasets and Protocol}
\label{sec:exp:data}

We deliberately avoid in-house authored evaluation slices for the
primary results: a method that beats a self-authored hard-negative
set has no reviewer-credible authority. All headline experiments
use established public benchmarks with externally published
baselines.

\paragraph{Primary benchmark: AgentDojo}~\citep{debenedetti2024agentdojo}.
The AgentDojo injection-task suite (97 user tasks $\times$ 27
injection tasks across four domains---banking, slack, travel,
workspace) is the closest public match to LSVJ-S's threat surface:
prompt-injection-mediated harm through legitimate tool calls.
Owner-Harm companion work~\citep{zhang2026ownerharm} reports the
current Nous gate at \textbf{14.8\%} TPR (4/27, 95\% CI:
5.9--32.5\%) on this slice under ground-truth deterministic
evaluation. This is the headline baseline LSVJ-S must beat.

\paragraph{Secondary benchmark: ShieldAgent-Bench}~\citep{chen2025shieldagent}.
The ICML 2025 ShieldAgent paper releases a multi-domain
rule-violation benchmark with $\approx 3$k traces. We adopt the
official splits and the published ShieldAgent baseline numbers as a
second comparison point; this is structurally important because
ShieldAgent itself is one of the closest related-work methods we
differentiate against in Section~\ref{sec:related}.

\paragraph{Reference / saturation slice: AgentHarm}~\citep{andriushchenko2025agentharm}.
The 176H + 176B AgentHarm slice is included in Appendix~C purely
as a no-regression check: the Owner-Harm companion gate already
saturates this slice at 100\% TPR / 0\% FPR; LSVJ-S simply needs
to maintain that.

\paragraph{Author-internal artifacts (NOT primary evidence).}
The repository also contains
\texttt{nous/data/owner\_harm\_heldout\_v3.json} (450 in-house
cases) and an early protocol sketch
\texttt{nous/docs/v3-HN-blind-protocol.md} for an in-house
blind-authored hard-negative slice. These are retained as design
and engineering artifacts only. \emph{No experimental claim in
this paper rests on them.} A pre-pivot version of the experiment
plan exists in repo history (commits up to \texttt{9f8c7b1}) and
is preserved for transparency, not used as evidence.

\paragraph{Statistical policy.}
All proportions are reported with Wilson 95\% CI; significance
tests at $\alpha = 0.05$; paired-McNemar is used only for the B4
decisive-subset analysis on AgentDojo's per-task per-injection grid.

\paragraph{Metrics.}
TPR/FPR with Wilson 95\% CI per category; paired-McNemar
$\chi^2$ with continuity correction; paired-Wilson 95\% CI on the
flip rate $b/(b+c)$; false-allow rate (B3 diagnostic); verbatim-match
rate (B2 diagnostic); parse-failure rate; decisive-primitive
distribution; per-call latency.

\subsection{B1 --- Main Comparison on Public Benchmarks}
\label{sec:exp:b1}

\paragraph{Configurations on AgentDojo.}
\begin{enumerate}[leftmargin=1.5em,nosep]
\item \textbf{No defense}: agent-pipeline only, no gate
(AgentDojo published baseline).
\item \textbf{Llama-Guard~3}~\citep{inan2023llamaguard}:
representative LLM-only safety classifier.
\item \textbf{Nous current}~\citep{zhang2026ownerharm}: L1+L2+L3(old
prompt)+L4 --- the 14.8\% baseline from the companion paper.
\item \textbf{LSVJ-S (ours)}: L1+L2+L3-LSVJ-S+L4;
\texttt{upgrade\_only=True}; majority-vote $k{=}3$; $N{=}5$
perturbations in b.4.
\end{enumerate}

\paragraph{Configurations on ShieldAgent-Bench.}
ShieldAgent-Bench~\citep{chen2025shieldagent} ships with its own
official baselines. We add LSVJ-S as a single new row and report
the same metrics defined by the original benchmark
(rule-compliance accuracy + utility preservation).

\paragraph{Success criterion (pre-registered).}
On AgentDojo: LSVJ-S TPR Wilson 95\% CI lower bound $>$ Nous-current
upper bound (i.e.\ statistically distinct improvement over the
14.8\% baseline). On ShieldAgent-Bench: LSVJ-S rule-compliance
$\ge$ ShieldAgent published number on at least one domain split.

\paragraph{Failure off-ramp (pre-registered).}
If LSVJ-S $\le$ Nous-current on AgentDojo statistically: pause and
pivot to PRM/RL-light or reasoning-model-native synthesis. We do
not reframe to hide a null result.

\placeholder{Table 1 (AgentDojo). TBD after eval. Columns:
configuration $\times$ injection-task category (TPR / FPR with
Wilson 95\% CI). Banking / slack / travel / workspace separately.}

\placeholder{Table 1b (ShieldAgent-Bench). TBD after eval.
ShieldAgent vs LSVJ-S on the official splits.}

\subsection{B2 --- Synthesis vs.\ Menu: Novelty Isolation}
\label{sec:exp:b2}
LSVJ-S full vs.\ SEAP (from B1) vs.\ \texttt{no\_synth\_6seed}
(Tier-2, Appendix~D). Key diagnostic: verbatim-match rate;
target $\le 30\%$; values above $50\%$ trigger the ``synthesis
collapsed to template completion'' finding. We also report distinct
rule bodies per 100 cases and the per-category TPR drop from full
LSVJ-S to SEAP. \placeholder{Table 2a TBD after W6.}

\subsection{B3 --- Stage + Class Necessity Ablation}
\label{sec:exp:b3}
Nine Tier-1 conditions, all must-run:
(1) LSVJ-S full;
(2) SEAP;
(3) \texttt{no\_sem\_gate} (b.4 disabled; measures false-allow on
neg-control, diagnostic only);
(4--6) \texttt{no\_class\_A}, \texttt{no\_class\_B},
\texttt{no\_class\_C} (each class disabled in turn);
(7) \texttt{KG-shuffled-all};
(8) \texttt{corrupt\_is\_inner\_circle\_only};
(9) \texttt{corrupt\_owner\_has\_directed\_only}.
Conditions 8 and 9 are paired with the B4 McNemar analysis.

\paragraph{Success criterion.}
\texttt{no\_sem\_gate} directionally higher false-allow on
neg-control; each \texttt{no\_class\_A/B/C} yields overall TPR drop
$\ge 3$\,pp. \placeholder{Table 2b TBD after W7.}

\subsection{B4 --- Per-Primitive Causal Attribution}
\label{sec:exp:b4}
This is the make-or-break test for the ``symbolic substrate
verifier'' framing: is the KG causally load-bearing, or is the
synthesized Datalog rule functionally a high-latency self-consistency
check?

\paragraph{Protocol.}
For each of two Tier-1 primitives
(\texttt{is\_inner\_circle}, \texttt{owner\_has\_directed}):
(i) run intact LSVJ-S on v3-HN; identify the decisive-subset via
\texttt{proof\_trace};
(ii) re-run with only that primitive's KG corrupted (within-type
shuffle);
(iii) compute paired McNemar with $b{=}$ intact-allow $\to$
corrupted-block; $c{=}$ intact-block $\to$ corrupted-allow; $\chi^2$
with continuity correction at $\alpha = 0.05$.

\paragraph{Effect-size floor (disclosed in main body).}
Expected decisive-subset size 20--30 per primitive; realistic
$b+c \approx 7\text{--}12$; McNemar significance at
$\alpha = 0.05$ requires $b/(b+c) \gtrsim 0.8$. Weaker effects are
reported as point estimate with paired-Wilson CI only, labeled
\emph{suggestive but below significance floor at current
decisive-subset size}.

\paragraph{Failure path.}
If both Tier-1 primitives produce null results, we test two Tier-2
primitives (\texttt{recipient\_owned\_by\_owner},
\texttt{action\_in\_precedent}); if all four null, we report the
honest null, pivot to JIT-policy-gate framing
(\S\ref{sec:disc:pivot}), and propose a KG-heavy obligation redesign
as future work.

\placeholder{Table 3 and Figure 3 TBD after W7.}

\subsection{B6 --- Adversarial Synthesis Stress Test}
\label{sec:exp:b6}
Empirical test of the Self-Justification Trap defense: 50
attacker-mode cases (hand-crafted) plus 30 v3-HN benign cases as
control. The proposer is given a sealed adversarial prompt
instructing it to bypass the gate by emitting a rule that passes
b.1--b.3 while remaining semantically vacuous.

\paragraph{Attack types.}
(i) tautology; (ii) permissive (rarely-enforcing primitives only);
(iii) collusion (crafted to pass perturbation surface checks);
(iv) multi-primitive OR-disjunction where only one primitive drives
discharge.

\paragraph{Success criterion.}
b.4 rejection rate on attacker-mode $\ge 90\%$; false-rejection on
benign control $\le 10\%$; per-attack-type rejection $\ge 75\%$
across all four types. \placeholder{Table 4 TBD.}

% ================================================================
\section{Discussion}
\label{sec:disc}

\subsection{Interpreting the Main Results}
\placeholder{Populated after B1--B4 results. If Claim~1 holds
(LSVJ-S $>$ SEAP with Wilson CI separation): per-decision runtime
synthesis adds TPR beyond an expert-authored policy of the same
vocabulary---the architectural delta PCAS cannot claim for
open-vocabulary cases. If Claim~2 locus (c) holds
(B4 paired-McNemar significant): the KG is causally load-bearing;
the synthesized Datalog rule is a proof obligation grounded in
system state.}

\subsection{Contingent Framing (Pre-Committed)}
\label{sec:disc:pivot}
Two framings are named before the results are observed:
\begin{description}[style=nextline]
\item[Primary framing] (if B4 shows a causal effect on $\ge 1$
Tier-1 primitive): \emph{LSVJ-S as symbolic substrate verifier.}
Per-decision Datalog synthesis grounded in live KG facts; the
compile-time gate ensures non-trivial commitment.
\item[Fallback framing] (if B4 null on all Tier-1 and Tier-2):
\emph{LSVJ-S as JIT-policy gate.} The LLM synthesizes a Datalog
rule as an explicit commitment object; b.4 filters self-justification
from genuine content-dependence. B6 becomes the primary evidence.
The ``symbolic substrate verifier'' phrasing is dropped. The method
thesis is unchanged; only the contribution-focus prose shifts.
\end{description}
This pivot is pre-registered
(\texttt{refine-logs/FINAL\_PROPOSAL.md}, Patch-3), not a post-hoc
save. Both framings are disclosed to the reader.

\subsection{Statistical Honesty: The Effect-Size Floor}
At the expected decisive-subset size (20--30 per primitive),
paired-McNemar significance at $\alpha = 0.05$ requires
$b/(b+c) \gtrsim 0.8$. This is disclosed in the main body, not hidden
in an appendix, with explicit explanation that the floor is a
function of sample size, not evidence of no effect. Weaker effects
($b/(b+c) \in [0.5, 0.8)$) are reported as point estimate plus
paired-Wilson CI with the label \emph{suggestive but below
significance floor}.

\subsection{Limitations}
We name, rather than bury, the workshop-tier limitations:
\begin{itemize}[leftmargin=1.5em,nosep]
\item \textbf{Self-Justification Trap.} Acknowledged; b.3 + b.4 are
the defenses; B6 is the empirical bound; b.4 does not catch all
collusion modes and we say so.
\item \textbf{Class C same-model oracle.} Procedural independence
only; representation-level independence is future work.
\item \textbf{KG-empty ablation deferred to Tier-2.} Workshop causal
attribution rests on two per-primitive corruptions plus
\texttt{KG-shuffled-all}; absolute KG absence is not tested at
workshop tier.
\item \textbf{v3-HN co-authorship.} The first author co-authors 50\%
of v3-HN and the SEAP menu. Pre-registration in
\texttt{nous/docs/no-synth-menu-plus.cozo} before any v3-HN
evaluation mitigates; independent human authoring is ICLR-tier
contingency.
\item \textbf{No soundness theorem.} Empirical evidence only (B6);
formal soundness analysis is future work.
\item \textbf{Workshop-tier novelty.} Self-Consistency
\citep{wang2022selfconsistency} covers b.4's mechanism. LSVJ-S
novelty is the combination applied to runtime safety-rule synthesis
plus decisive-primitive attribution. We do not claim main-track
status.
\end{itemize}

% ================================================================
\section{Conclusion}
\label{sec:conclusion}
LSVJ-S instantiates the Synthesis--Proof-Obligation (SPO) protocol:
the LLM synthesizes a per-decision Datalog rule as its committed
proof of trust, and a four-stage compile-time gate---parse,
type-check, syntactic non-triviality, and perturbation-sensitive
\emph{and} has-decisive-primitive---admits the rule only when all
stages pass. Stage b.4 is grounded in Self-Consistency
\citep{wang2022selfconsistency}; its combination with typed Datalog
synthesis and decisive-primitive causal attribution is our novel
contribution in the runtime agent-safety context.
\placeholder{Headline empirical summary TBD after W6--W7.}
Named limitations include Class~C same-model procedural
independence, the absence of a formal soundness theorem, and the
workshop scope. Future work: different-model Class~C, formal
soundness, an external zero-shot benchmark such as
ShieldAgent-Bench.

% ================================================================
\section*{Author Contributions}
\textbf{Dongcheng Zhang} designed the SPO protocol and the 4-stage
gate, authored the M0 implementation, led the R002 pilot, authored
the 50\% of v3-HN that is not blind-slice, and drafted the
manuscript. \textbf{Yiqing Jiang} contributed the knowledge-graph
primitive design and provided critical review of the ontology,
method, and experimental protocol.

\section*{Acknowledgments}
We thank GPT-5.4 (via Codex), Opus~4.7, and Gemini~3.1~Pro (via
Gemini CLI) for independent critical audits across four refinement
rounds.

% ================================================================
\bibliographystyle{plainnat}
\bibliography{references}

\end{document}
